\section{A global residue functional and its interval certificate}
\label{boundary:app:residue:rr}

We work at phase zero, with $P_\alpha=Q_0-\alpha J$.
For real $0<\alpha<g$ put
\[
 a=\sqrt{1-2\alpha^2},\quad t=\arctan a,\quad
 \nu=\tfrac12+t/\pi,\quad c=(2\cos t)^{-1},\quad q=ce^{-it}.
\]
For $p>0$ and real $z<1$ define
\begin{align}
 \ell_p(z)&=\frac1\pi\int_0^1\frac{x^{p-1}}{1-zx}\,dx,
 \nonumber\\
 A_\nu(s)&=-\int_0^1\frac{x^{-\nu}}{1+s-x}\,dx
   =-\frac\pi{1+s}\ell_{1-\nu}\!\left(\frac1{1+s}\right),
 \qquad B_\nu(s)=-A_\nu(1-s).
 \label{boundary:eq:residue-scalar-kernels:rr}
\end{align}
Let
\[
 k=(1,(\alpha/c)e^{it},i(\alpha/c)e^{2it},ie^{3it})
\]
be a row, and define
\begin{align}
 j_1(s)&=k_2A_\nu(s)+\pi[-\ell_\nu(s)
                  +k_3\ell_\nu(1-s)+k_4\ell_\nu(-s)],\nonumber\\
 j_2(s)&=k_3B_\nu(s)+\pi[-\ell_\nu(s-1)
                  -k_2\ell_\nu(s)+k_4\ell_\nu(1-s)],\nonumber\\
 p_\alpha&=\frac{-i\Gamma(\nu)e^{-i(\pi/4+3t/2)}}
                  {8\pi^2ac(2\pi)^\nu},\qquad
 F_\alpha(u)=p_\alpha\int_0^1
                [j_1(s)u(s)+j_2(s)u(s-1)]\,ds.
 \label{boundary:eq:global-residue-functional:rr}
\end{align}
The scalar $p_\alpha$ here is a prefactor, distinct from the pencil
$P_\alpha$.  The two row kernels have only logarithmic endpoint
singularities, so $F_\alpha$ is a bounded complex-linear functional
on $L^2(-1,1)$.

\begin{lemma}[The global residue row]
\label{boundary:lem:global-residue:rr}
If $P_\alpha u=0$, then $F_\alpha(u)$ is the inward physical tail
amplitude of $V_0u$, with its induced physical normalization.
Moreover, in the operator norm of functionals,
\begin{equation}\label{boundary:eq:uniform-residue-bounds:rr}
 \|F_\alpha\|<\frac{13}{100},\qquad
 \|\partial_\alpha F_\alpha\|<1,
 \qquad \frac25\le\alpha\le\frac{11}{25}.
\end{equation}
\end{lemma}

\subsection{Derivation and the auxiliary cutoff limit}

The right null column of the endpoint multiplier at $\omega=-i\beta$,
with first component one, is
\[
 v=(1,(\alpha/c)e^{it},-i(\alpha/c)e^{2it},-ie^{3it})^T.
\]
Direct multiplication in \eqref{v12:eq:endpoint-symbol} shows that
$k$ is its left null row, without conjugation.  Its adjugate has
$(1,1)$ entry $-q^3$ and
\[
 \partial_\omega\det m_{\alpha,0}(-i\beta)
       =i4\pi\alpha^2ac^2.
\]
For a smooth endpoint cutoff, the residue formula in the proof of
Theorem~\ref{v12:thm:boundary-tail} therefore gives the endpoint vector
\begin{equation}\label{boundary:eq:source-residue-row:rr}
 V=\frac{q^3}{4\pi\alpha^2ac^2}\,vN,\qquad
 N=\int_0^\infty x^{-\nu}k(-\Delta_eu)(x)\,dx.
\end{equation}
Indeed the factor $-i\sqrt{2\pi}$ in that residue formula cancels
the Mellin normalization and the displayed determinant derivative.
The positive zero-carrier Gamma factor for $v$ is
$2\cos(2t)e^{-i\pi/4+3it/2}$, and
$\alpha^2=2c^2\cos(2t)$.  Thus its physical amplitude is
\[
 \frac{\Gamma(\nu)e^{-i(\pi/4+3t/2)}}
      {4\pi ac(2\pi)^\nu}\,N.
\]

This calculation uses only $EP_\alpha=M_\alpha E+\Delta_e$ and
$\chi=1$ near zero; the auxiliary cutoffs need not form an isometric
partition. For two such cutoffs the difference localizes away from
zero. On a kernel vector its Mellin transform is regular at $-i\beta$
and $km_{\alpha,0}(-i\beta)=0$, so the source pairing is unchanged.
No cutoff independence on arbitrary $u$ is asserted.

Take smooth cutoffs equal to one up to $1-2\varepsilon$ and zero
after $1-\varepsilon$, and let $\varepsilon\downarrow0$.
Their source-pairing rows converge in $L^2$ to the row for
$\chi=\mathbf1_{(0,1)}$.  Here are the required operator bounds.
In $EQ_0$, move $Q_0$ onto the cutoff difference multiplied by
$x^{-\nu}k$; this weight is bounded near one and $Q_0$ is bounded
on $L^2$.  In $Q_eE$, move $Q_e$ onto $x^{-\nu}k$ and restrict the
result to a neighborhood of one.  That result is bounded: the input
near zero is integrable since $\nu<1$, its tail is integrable against
the Cauchy denominator, and its local principal value is the transform
of a smooth power.  Both error rows tend to zero with the $L^2$ norm
of the cutoff difference.  These statements may first be proved for
smooth inputs and then extended by density.  They require no uniform
derivative bound on the steep cutoffs.

For this limiting cutoff, the kernel in
\eqref{v12:eq:endpoint-defect-kernel}, before its factor $i/(2\pi)$,
has on $0<x<1$ only the entries
$(2,1)=(x-1-s)^{-1}$ and $(3,2)=(2-x-s)^{-1}$.
On $x>1$ it is
\[
 \begin{pmatrix}
 -(x-s)^{-1}&-(x+1-s)^{-1}\\
 0&-(x-s)^{-1}\\
 (x+s-1)^{-1}&0\\
 (x+s)^{-1}&(x+s-1)^{-1}
 \end{pmatrix}.
\]
Integration against $x^{-\nu}k$ gives exactly the two rows in
\eqref{boundary:eq:global-residue-functional:rr}, with
$N=-i(2\pi)^{-1}\int_0^1(j_1u_1+j_2u_2)$.
This proves the first assertion of
Lemma~\ref{boundary:lem:global-residue:rr}, including its normalization.

\subsection{Uniform row and derivative bounds}

All decimals in the following estimates are finite rational numbers.
For $2/5\le\alpha\le11/25$, elementary differentiation gives
\begin{gather*}
 a>0.78,\quad c>0.63,\quad 0.7<\nu<0.75,\quad
 \kappa:=\alpha/c<0.7,\\
 |t'|<0.71,\quad |\nu'|<0.23,\quad |c'|<0.35,\quad
 |\kappa'|<2,\quad |k_2'|<2.5,\quad |k_3'|<3,\quad |k_4'|<2.14.
\end{gather*}
For example,
$t'=-2\alpha/[a(1+a^2)]$, $c'=-\alpha/(2c)$, and
$\kappa'=1/c+\alpha^2/(2c^3)$.  In bounding $\nu'$ below we also
retain the sharper bound $|\nu'|<0.71/\pi$.

The power series
$\pi\ell_p(z)=\sum_{n\ge0}z^n/(n+p)$ for $0\le z<1$, and the
positive defining integrals for $z\le0$, imply
\begin{gather*}
 \pi\ell_\nu(s)\le\nu^{-1}-\log(1-s),\qquad
 \pi\ell_\nu(-s),\ \pi\ell_\nu(s-1)\le\nu^{-1},\\
 |A_\nu(s)|\le(1-\nu)^{-1}+\log2-\log s,\qquad
 B_\nu(s)\le(1-\nu)^{-1}+\log2-\log(1-s).
\end{gather*}
For the last bounds apply the series at $z=(1+s)^{-1}$ and its
reflection.  Differentiation in $\nu$ gives the absolute bounds
\[
 \nu^{-2}+\pi^2/6<15/4,\qquad \nu^{-2}<21/10,
 \qquad (1-\nu)^{-2}+\pi^2/6<18,
\]
respectively for a positive-argument $\pi\ell$, a nonpositive-argument
$\pi\ell$, and $A,B$.  The series bounds use
$\sum_{n\ge1}(n+p)^{-2}\le\pi^2/6$.

With $L_0=-\log s$ and $L_1=-\log(1-s)$, these estimates yield
\begin{align}
 |j_1|&\le7.2+1.4L_0+L_1,&
 |j_2|&\le7.2+L_0+1.4L_1,\nonumber\\
 |j_1'|&\le24+5.5L_0,&
 |j_2'|&\le26+2.2L_0+5.5L_1.
 \label{boundary:eq:logarithmic-row-bounds:rr}
\end{align}
For example the first undifferentiated constant is at most
$0.7(4.7)+2.7(10/7)<7.2$.  The first differentiated constant is at most
\[
 2.5(4.7)+(3+2.14)\frac{10}{7}
 +0.23\left(0.7(18)+1.7\frac{15}{4}+\frac{21}{10}\right)<24.
\]
Replacing $2.5(4.7)+3(10/7)$ by $3(4.7)+2.5(10/7)$ bounds the
second by 26.  The logarithmic coefficients in
\eqref{boundary:eq:logarithmic-row-bounds:rr} are respectively the
sums of the corresponding $|k_j|$ or $|k_j'|$.

The scalar prefactor satisfies
\begin{equation}\label{boundary:eq:prefactor-bounds:rr}
 |p_\alpha|<0.0093,\qquad |p_\alpha'|<4|p_\alpha|.
\end{equation}
For the first inequality use
\[
 |p_\alpha|\le
 \frac{\Gamma(0.7)}{(2\pi)^{0.7}\,8\pi^2(0.78)(0.63)}<0.0093.
\]
The digamma $\psi_0=\Gamma'/\Gamma$ is increasing, since
$\psi_0'(x)=\sum_{n\ge0}(n+x)^{-2}>0$.
The scalar bounds $\psi_0(0.7)>-1.3$ and $\psi_0(0.75)<0$ show that
$\Gamma$ decreases and $|\psi_0|<1.3$ on $[0.7,0.75]$.
Consequently
\[
 |p_\alpha'/p_\alpha|
 \le\frac{0.71}{\pi}(1.3+\log(2\pi))
       +\frac{0.88}{0.78^2}+\frac{0.35}{0.63}+\tfrac32(0.71)<4.
\]
All scalar inequalities here are also checked with outward-rounded
arithmetic by the accompanying verifier.

Finally, $\int L_j=1$, $\int L_j^2=2$, and
$\int L_0L_1=2-\pi^2/6<0.36$.
Squaring \eqref{boundary:eq:logarithmic-row-bounds:rr} and integrating
gives bounds $186.656$ for $\|(j_1,j_2)\|_2^2$ and $2055.792$ for
$\|(j_1',j_2')\|_2^2$.  Therefore
\[
 \|F_\alpha\|<0.0093\sqrt{186.656}<0.13,\qquad
 \|\partial_\alpha F_\alpha\|
 <4(0.13)+0.0093\sqrt{2055.792}<1.
\]
The same logarithmic envelopes justify differentiation in operator
norm.  This finishes Lemma~\ref{boundary:lem:global-residue:rr}.

\subsection{Exact integration of the rational witness}

Define primitives
\begin{align}
 E_\nu(z)&=\frac{z\ell_\nu(z)+\pi^{-1}\log(1-z)}{1-\nu},
 &E_\nu'&=\ell_\nu,\nonumber\\
 D_\nu(s)&=\frac{(1+s)A_\nu(s)-\log s}{1-\nu},
 &D_\nu'&=A_\nu.
 \label{boundary:eq:residue-primitives:rr}
\end{align}
Their continuous endpoint values are
\[
 E_\nu(0)=0,\quad
 E_\nu(1)=-\frac{\psi_0(\nu)+\gamma_E}{\pi(1-\nu)},\quad
 D_\nu(0)=\frac{\psi_0(1-\nu)+\gamma_E}{1-\nu},
\]
where $\gamma_E$ is Euler's constant.
The derivative identities follow by integration by parts in the
defining integrals; the endpoint values follow from
$\pi\ell_p(z)+\log(1-z)\to-\psi_0(p)-\gamma_E$ as $z\uparrow1$.
Primitives of $j_1,j_2$ are consequently
\begin{align}
 J_1(s)&=k_2D_\nu(s)+\pi[-E_\nu(s)-k_3E_\nu(1-s)-k_4E_\nu(-s)],
 \nonumber\\
 J_2(s)&=k_3D_\nu(1-s)+\pi[-E_\nu(s-1)-k_2E_\nu(s)-k_4E_\nu(1-s)].
 \label{boundary:eq:residue-row-primitives:rr}
\end{align}

Use the same exact witness $u_0$ and parameter $a_0$ as in
\eqref{rr:eq:positive-trial}.
Integrate \eqref{boundary:eq:residue-row-primitives:rr} over each
cell and multiply by the witness's constant value there.  Its stored
coefficients multiply normalized indicators, so this value is the
coefficient divided by the square root of the exact cell width.
The first 256 coefficients belong to the negative half and the last
256 to the positive half; they pair with $j_2$ and $j_1$, respectively.

The accompanying interval verifier evaluates these endpoint
differences using
$\ell_p(z)={}_2F_1(1,p;p+1;z)/(\pi p)$, Gamma and digamma functions
with outward-rounded Arb arithmetic \cite{Johansson2017}.
The exact parameter relations $b-c=-1$ and $a+b-c=0$ are supplied
when evaluating the zero-balanced hypergeometric function. At both
192 and 256 bits it verifies
\begin{equation}\label{boundary:eq:exact-residue-anchor:rr}
 |F_{a_0}(u_0)|>0.069,
 \qquad \|u_0\|<1.001,
 \qquad 0.41635<a_0<0.41637,
\end{equation}
as well as the scalar bounds used above.  No quadrature, eigensolver,
or saved floating-point result enters these checks.
The existing residual certificate gives
\begin{equation}\label{boundary:eq:frequency-residual-bound:rr}
 \|P_{a_0}u_0\|<0.00622:
 \qquad U_{u_0}(a_0)<0.000038641<(0.00622)^2.
\end{equation}

\subsection{The complete physical projection estimate}

\begin{proof}[Nonvanishing in Theorem~\ref{boundary:thm:certified-residue:rr}]
By the proved isolation clause, $\alpha_0\in(0.405,0.428)$ is the
only spectral point in $[0.3565,0.4765]$. Let $\Pi$ project onto it.
Equation~\eqref{boundary:eq:exact-residue-anchor:rr} puts every other
spectral point at distance at least $0.049$ from $a_0$.

Put $f_0=V_0u_0$.  The identities
$B_0=V_0JV_0^*$, $V_0^*V_0=Q_0$, and $\|V_0\|\le1$ give
\[
 (B_0-a_0)f_0=V_0JP_{a_0}u_0,\qquad
 \|(I-\Pi)f_0\|<\frac{0.00622}{0.049}.
\]
For $\alpha\in(0.405,0.428)$ define the bounded physical functional
\[
 \mathcal L_\alpha f=\alpha^{-1}F_\alpha(JV_0^*f),\qquad
 \|\mathcal L_\alpha\|<\frac{0.13}{0.405}.
\]
On the $\alpha$ eigenspace it equals the physical tail amplitude:
$u=\alpha^{-1}JV_0^*f$ satisfies $P_\alpha u=0$ and $V_0u=f$.
For the trial vector, the exact identity $JQ_0=\alpha I+JP_\alpha$
instead gives
\[
 \mathcal L_\alpha f_0
    =F_\alpha(u_0)+\alpha^{-1}F_\alpha(JP_\alpha u_0).
\]
Thus no trial value is identified with an eigenvector value.

Now $|\alpha_0-a_0|<0.012$, and
\eqref{boundary:eq:uniform-residue-bounds:rr} yields
\[
 |F_{\alpha_0}(u_0)-F_{a_0}(u_0)|<0.012(1.001),\qquad
 \|P_{\alpha_0}u_0\|<0.00622+0.012(1.001).
\]
Combining these bounds with the full complementary projection error,
\begin{align*}
 |\mathcal L_{\alpha_0}\Pi f_0|
 &>0.069-0.012(1.001)\\
 &\quad-\frac{0.13}{0.405}
       \left(0.00622+0.012(1.001)+\frac{0.00622}{0.049}\right)
 >0.0103.
\end{align*}
The final inequality is rational and is checked independently by the
same verifier.  Since $\|\Pi f_0\|\le\|f_0\|\le\|u_0\|<1.001$,
normalizing the nonzero eigenfunction $\Pi f_0$ gives $|\eta|>0.01$.
Simplicity makes this magnitude independent of the phase of the
normalized eigenfunction.  The asserted parity consequence now
follows from Corollary~\ref{v13:cor:leading-parity-splitting}.
\end{proof}

